\documentclass[11pt,a4paper]{article}

\usepackage[utf8]{inputenc}
\usepackage[T1]{fontenc}
\usepackage{amsmath,amssymb,mathtools}
\usepackage{xcolor}
\usepackage{mdframed}
\usepackage{geometry}
\usepackage{microtype}

\geometry{
  margin=2.4cm
}

\setlength{\parindent}{0pt}
\setlength{\parskip}{0.7em}

% Couleurs imposees
\definecolor{bleutheoreme}{RGB}{31,78,121}
\definecolor{fondbleu}{RGB}{234,242,250}

\definecolor{orangelemme}{RGB}{210,105,30}
\definecolor{fondorange}{RGB}{255,244,230}

\definecolor{vertpreuve}{RGB}{35,130,75}
\definecolor{fondvert}{RGB}{235,248,239}

\definecolor{violetremarque}{RGB}{112,48,160}
\definecolor{fondviolet}{RGB}{246,238,252}

\definecolor{rougeresultat}{RGB}{180,30,45}
\definecolor{fondrouge}{RGB}{255,235,238}

\newmdenv[
  linecolor=bleutheoreme,
  backgroundcolor=fondbleu,
  linewidth=1.2pt,
  roundcorner=4pt,
  skipabove=10pt,
  skipbelow=10pt
]{theoremebox}

\newmdenv[
  linecolor=orangelemme,
  backgroundcolor=fondorange,
  linewidth=1.2pt,
  roundcorner=4pt,
  skipabove=10pt,
  skipbelow=10pt
]{lemmebox}

\newmdenv[
  linecolor=vertpreuve,
  backgroundcolor=fondvert,
  linewidth=1.2pt,
  roundcorner=4pt,
  skipabove=10pt,
  skipbelow=10pt
]{preuvebox}

\newmdenv[
  linecolor=violetremarque,
  backgroundcolor=fondviolet,
  linewidth=1.2pt,
  roundcorner=4pt,
  skipabove=10pt,
  skipbelow=10pt
]{remarquebox}

\newmdenv[
  linecolor=rougeresultat,
  backgroundcolor=fondrouge,
  linewidth=1.5pt,
  roundcorner=4pt,
  skipabove=12pt,
  skipbelow=12pt
]{resultatbox}

\newcommand{\N}{\mathbb{N}}
\newcommand{\R}{\mathbb{R}}

\title{\textbf{Somme de puissances dont l'exposant cro\^{i}t avec \(n\)}}
\author{R\'edaction rigoureuse pour Mostafa}
\date{12 juillet 2026}

\begin{document}

\maketitle

\begin{theoremebox}
\textbf{\'Enonc\'e.}
Soit \(\alpha>0\). D\'emontrer que, lorsque \(n\) tend vers
\(+\infty\),
\[
1^{\alpha n}+2^{\alpha n}+\cdots+n^{\alpha n}
\sim
\frac{n^{\alpha n}}{1-e^{-\alpha}}.
\]
\end{theoremebox}

\begin{lemmebox}
\textbf{Lemme 1.}
Pour tout \(x\in[0,1)\),
\[
\log(1-x)\leq -x.
\]
En particulier, pour \(0\leq k<n\),
\[
0<
\left(1-\frac{k}{n}\right)^{\alpha n}
\leq e^{-\alpha k}.
\]
\end{lemmebox}

\begin{preuvebox}
\textbf{Preuve du lemme 1.}
La fonction
\[
h(x)=\log(1-x)+x
\]
est d\'erivable sur \([0,1)\), et
\[
h'(x)=-\frac{1}{1-x}+1=-\frac{x}{1-x}\leq 0.
\]
Comme \(h(0)=0\), on obtient \(h(x)\leq0\), donc
\[
\log(1-x)\leq-x.
\]
En prenant \(x=k/n\) et en multipliant par \(\alpha n>0\),
\[
\alpha n\log\left(1-\frac{k}{n}\right)
\leq-\alpha k.
\]
L'exponentiation conserve l'in\'egalit\'e, ce qui donne
\[
\left(1-\frac{k}{n}\right)^{\alpha n}
\leq e^{-\alpha k}.
\]
\end{preuvebox}

\begin{lemmebox}
\textbf{Lemme 2.}
Pour tout \(x\in[0,1/2]\),
\[
0\leq-\log(1-x)-x\leq x^2.
\]
\end{lemmebox}

\begin{preuvebox}
\textbf{Preuve du lemme 2.}
Pour \(x\in[0,1)\), le d\'eveloppement en s\'erie enti\`ere de
\(-\log(1-x)\) donne
\[
-\log(1-x)-x
=\sum_{m=2}^{\infty}\frac{x^m}{m}.
\]
Tous les termes sont positifs. De plus, \(1/m\leq1/2\) pour
\(m\geq2\). Par cons\'equent,
\[
0\leq-\log(1-x)-x
\leq
\frac12\sum_{m=2}^{\infty}x^m
=
\frac{x^2}{2(1-x)}.
\]
Si \(x\leq1/2\), alors \(2(1-x)\geq1\), d'o\`u
\[
-\log(1-x)-x\leq x^2.
\]
\end{preuvebox}

\begin{preuvebox}
\textbf{Preuve principale.}

Posons
\[
S_n=\sum_{j=1}^{n}j^{\alpha n}.
\]
Apr\`es division par \(n^{\alpha n}\), puis changement d'indice
\(k=n-j\), nous obtenons
\[
\frac{S_n}{n^{\alpha n}}
=
\sum_{j=1}^{n}\left(\frac{j}{n}\right)^{\alpha n}
=
\sum_{k=0}^{n-1}
\left(1-\frac{k}{n}\right)^{\alpha n}.
\]

Nous allons comparer cette somme \`a la s\'erie g\'eom\'etrique
\[
\sum_{k=0}^{\infty}e^{-\alpha k}
=
\frac{1}{1-e^{-\alpha}}.
\]

Pour \(1\leq k<n\), le lemme 1 donne
\[
0<
\left(1-\frac{k}{n}\right)^{\alpha n}
\leq e^{-\alpha k}.
\]
Ainsi,
\[
\left|
\left(1-\frac{k}{n}\right)^{\alpha n}
-e^{-\alpha k}
\right|
=
e^{-\alpha k}
-\left(1-\frac{k}{n}\right)^{\alpha n}.
\]

Conform\'ement \`a l'indication, nous scindons l'erreur \`a
l'indice \(\lfloor n/2\rfloor\). Posons
\[
E_n^{(1)}
=
\sum_{k=1}^{\lfloor n/2\rfloor}
\left|
\left(1-\frac{k}{n}\right)^{\alpha n}
-e^{-\alpha k}
\right|
\]
et
\[
E_n^{(2)}
=
\sum_{k=\lfloor n/2\rfloor+1}^{n-1}
\left|
\left(1-\frac{k}{n}\right)^{\alpha n}
-e^{-\alpha k}
\right|.
\]

Pour \(1\leq k\leq\lfloor n/2\rfloor\), posons
\[
\delta_{k,n}
=
-\log\left(1-\frac{k}{n}\right)-\frac{k}{n}.
\]
Le lemme 2 donne
\[
0\leq\delta_{k,n}\leq\frac{k^2}{n^2}.
\]
Nous avons
\[
\left(1-\frac{k}{n}\right)^{\alpha n}
=
e^{-\alpha k}e^{-\alpha n\delta_{k,n}}.
\]
L'in\'egalit\'e
\[
0\leq1-e^{-u}\leq u,
\qquad u\geq0,
\]
entra\^{i}ne
\[
\begin{aligned}
0
&\leq
e^{-\alpha k}
-\left(1-\frac{k}{n}\right)^{\alpha n}\\
&=
e^{-\alpha k}
\left(1-e^{-\alpha n\delta_{k,n}}\right)\\
&\leq
\alpha n e^{-\alpha k}\delta_{k,n}\\
&\leq
\frac{\alpha k^2}{n}e^{-\alpha k}.
\end{aligned}
\]
Par cons\'equent,
\[
E_n^{(1)}
\leq
\frac{\alpha}{n}
\sum_{k=1}^{\infty}k^2e^{-\alpha k}.
\]
La s\'erie
\[
\sum_{k=1}^{\infty}k^2e^{-\alpha k}
\]
converge, car \(\alpha>0\). Il s'ensuit que
\[
E_n^{(1)}\longrightarrow0.
\]

Pour la seconde partie, le lemme 1 donne
\[
E_n^{(2)}
\leq
\sum_{k=\lfloor n/2\rfloor+1}^{n-1}e^{-\alpha k}
\leq
\sum_{k=\lfloor n/2\rfloor+1}^{\infty}e^{-\alpha k}.
\]
La formule de sommation d'une s\'erie g\'eom\'etrique donne
\[
E_n^{(2)}
\leq
\frac{
 e^{-\alpha(\lfloor n/2\rfloor+1)}
}{
 1-e^{-\alpha}
}
\longrightarrow0.
\]

Enfin,
\[
\sum_{k=n}^{\infty}e^{-\alpha k}
=
\frac{e^{-\alpha n}}{1-e^{-\alpha}}
\longrightarrow0.
\]
Nous en d\'eduisons
\[
\begin{aligned}
&
\left|
\sum_{k=0}^{n-1}
\left(1-\frac{k}{n}\right)^{\alpha n}
-
\sum_{k=0}^{\infty}e^{-\alpha k}
\right|\\
&\hspace{2cm}
\leq
E_n^{(1)}+E_n^{(2)}
+\sum_{k=n}^{\infty}e^{-\alpha k}
\longrightarrow0.
\end{aligned}
\]
Ainsi,
\[
\frac{S_n}{n^{\alpha n}}
\longrightarrow
\sum_{k=0}^{\infty}e^{-\alpha k}
=
\frac{1}{1-e^{-\alpha}}.
\]
\end{preuvebox}

\begin{resultatbox}
\textbf{R\'esultat final.}
Pour tout \(\alpha>0\),
\[
\boxed{
1^{\alpha n}+2^{\alpha n}+\cdots+n^{\alpha n}
\sim
\frac{n^{\alpha n}}{1-e^{-\alpha}}
}
\qquad(n\longrightarrow+\infty).
\]
\end{resultatbox}

\begin{remarquebox}
\textbf{Remarque.}
La constante
\[
\frac{1}{1-e^{-\alpha}}
\]
provient des termes situ\'es \`a distance born\'ee de \(n\).
Pour chaque entier fix\'e \(k\geq0\),
\[
\left(\frac{n-k}{n}\right)^{\alpha n}
\longrightarrow e^{-\alpha k}.
\]
Les termes \(n^{\alpha n},(n-1)^{\alpha n},(n-2)^{\alpha n}\)
produisent donc asymptotiquement les contributions respectives
\[
1,\quad e^{-\alpha},\quad e^{-2\alpha},
\]
apr\`es normalisation par \(n^{\alpha n}\).
\end{remarquebox}

\begin{remarquebox}
\textbf{Classification.}
Cet exercice rel\`eve du th\`eme \textbf{Analyse}.
\end{remarquebox}

\end{document}